\cotefile{157-3}
\scanpage{192}\mspage{141}
et les flèches correspondantes

\diagram{\(\left\{\begin{aligned}
\varphi_Q &: \unsure{\mathscr X|\mathcal X}_{\deleted{\unsure{Q}}Q}=\left(\int \begin{tikzcd}[ampersand replacement=\&,row sep=0pt,column sep=1.2em,cramped] \& X_1 \\ X_0\arrow[ur,dash]\arrow[dr] \& \\ \& X_2 \end{tikzcd}\right)\longrightarrow Y\\
\varphi_{Q^\circ} &: \mathscr X^{\deleted{\unsure{Q}}}_{Q^\circ}=\left(\int \begin{tikzcd}[ampersand replacement=\&,row sep=0pt,column sep=1.2em,cramped] \& X_1^\circ \\ X_0^\circ\arrow[ur,dash]\arrow[dr,dash] \& \\ \& X_2^\circ \end{tikzcd}\right)\longrightarrow Y^\circ .
\end{aligned}\right.\)}

Enfin, considérons aussi les flèches $\widetilde\varphi_Q,\inserted{\unsure{\overline\varphi_{Q^\circ}|\widetilde\varphi_{Q^\circ}}}$ canoniques, relatives à l'intégrale propre $\int_{\mathrm{pr}}$ (au lieu de l'intégrale lisse $\int=\int_{\unsure{\mathrm{lis}|\mathrm{lic}}}$)

\diagram{\(\left\{\begin{aligned}
\widetilde\varphi_Q &\unsure{\mathrel{:}|\mathrel{=}} \widetilde{\mathscr X}_Q=\left(\int_{\mathrm{pr}} \begin{tikzcd}[ampersand replacement=\&,row sep=0pt,column sep=1.2em,cramped] \& X_1 \\ X_0\arrow[ur,dash]\arrow[dr] \& \\ \& X_2 \end{tikzcd}\right)\longrightarrow Y\\
\deleted{\text{\unsure{associées}}}\quad\deleted{\widetilde\varphi_{Q^\circ}:}&\deleted{\left(\int_{\mathrm{pr}} \begin{tikzcd}[ampersand replacement=\&,row sep=0pt,column sep=1.2em,cramped] \& X_1^\circ \\ X_0^\circ\arrow[ur,dash]\arrow[dr,dash] \& \\ \& X_2^\circ \end{tikzcd}\right)\longrightarrow \unsure{Y^\circ}}\\
\widetilde\varphi_{Q^\circ} &: \widetilde{\mathscr X}_{Q^\circ}=\left(\int \begin{tikzcd}[ampersand replacement=\&,row sep=0pt,column sep=1.2em,cramped] \& X_1^\circ \\ X_0^\circ\arrow[ur]\arrow[dr,dash] \& \\ \& X_2^\circ \end{tikzcd}\right)\longrightarrow Y^\circ .
\end{aligned}\right.\)}

On \unsure{a} alors ce qui suit :

\[
\text{\textcircled{\unsure{a}}}\quad \varphi_{Q^\circ}\simeq(\widetilde\varphi_Q)^\circ,\quad \varphi_Q\simeq(\widetilde\varphi_{Q^\circ})^\circ
\]

(d'où résulte déjà que $\left\{\begin{aligned}\varphi_Q\in W&\Longleftrightarrow\widetilde\varphi_{Q^\circ}\in W\\\varphi_{Q^\circ}\in W&\Longleftrightarrow\widetilde\varphi_Q\in W\end{aligned}\right.$ \unsure{\deleted{)}} ${}^{\star}$
% slanted
\margin{et que les flèches $\varphi_Q$ et $\widetilde\varphi_{Q^\circ}$ d'une part, $\varphi_{Q^\circ}$ et $\widetilde\varphi_Q$ d'autre part, sont \unsure{isomorphes} comme flèches dans $\mathrm{Hot}_W$.}).

\textcircled{b} Les relations

\[
\varphi_Q\in W,\quad \varphi_{Q^\circ}\in W,\quad \widetilde\varphi_Q\in W,\quad \widetilde\varphi_{Q^\circ}\in W
\]

sont équivalentes, en d'autres termes : pour que le carré soit $W$-homotopiquement cocartésien (version intégrale lisse), il f. et il suffit que le carré opposé $Q^\circ$ le soit, et ces conditions équivalent aussi aux conditions correspondantes en termes des intégrales propres (au lieu des intégrales lisses).

\textcircled{c} Les quatre flèches $\varphi_Q,\varphi_{Q^\circ},\widetilde\varphi_Q,\widetilde\varphi_{Q^\circ}$, \inserted{\unsure{de $\mathrm{Cat}$|de $\mathrm{Cat}$,} en tant que flèches} dans $\mathrm{Hot}_W$, sont isomorphes dans \uline{$\mathrm{Fl}$}$(\mathrm{Hot}_W)$. Plus précisément, en tant que flèches de $\mathrm{Cat}$,
% archive sheet number: 192
\scanpage{193}\mspage{142}
elles définissent des objets isomorphes de $\mathrm{HOT}_W(\Delta^1)$.

\textcircled{d} Plus précisément \unsure{encore|même}, les quatre diagrammes dans $\mathrm{Cat}$, de type $\Gamma^\circ$, où

\diagram{\(\Gamma={}\begin{tikzcd}[ampersand replacement=\&,row sep=.4em,column sep=1.2em]
\& 1\arrow[dl,"\alpha_1"'] \& \& \\
0 \& \& 3\arrow[ul,"i'_1"']\arrow[dl,"i'_2"] \& 4\arrow[l,"\varphi"']\\
\& 2\arrow[ul,"\alpha_2"] \& \&
\end{tikzcd},\quad\alpha_1i'_1\varphi=\alpha_2i'_2\varphi,\)}

dans lesquels s'insèrent les quatre flèches qu'on vient d'envisager, savoir

\diagram{\begin{tikzcd}[ampersand replacement=\&,row sep=.3em,column sep=1.4em]
\& \& X_1\arrow[dr,"i_1"] \& \& \\
\widehat Q={} \& X_0\arrow[ur,"\alpha_1"]\arrow[dr,dash,"\alpha_2"']\arrow[rr,dotted] \& \& \unsure{\mathscr X|\mathcal X}_Q\arrow[r,"\varphi_Q"] \& Y\\
\& \& X_2\arrow[ur,"i_2"'] \& \&
\end{tikzcd}}

% slanted
\margin{\uline{NB} On devrait \deleted{\unsure{et}} encore compléter ces diagrammes \deleted{par les} \unsure{(bien comm.)} par les flèches en pointillés\ldots}

\diagram{\begin{tikzcd}[ampersand replacement=\&,row sep=.3em,column sep=1.4em]
\& \& X_1^\circ\arrow[dr,"i'_1"] \& \& \\
\widehat{Q^\circ}_{\deleted{\unsure{\circ}}}={} \& X_0^\circ\arrow[ur,"\alpha_1^\circ"]\arrow[dr,"\alpha_2^\circ"']\arrow[rr,dotted] \& \& \mathscr X_{Q^\circ}\arrow[r] \& Y^\circ\\
\& \& X_2^\circ\arrow[ur,"i'_2"'] \& \&
\end{tikzcd}}

\diagram{\begin{tikzcd}[ampersand replacement=\&,row sep=.3em,column sep=1.4em]
\& \& X_1\arrow[dr] \& \& \\
\widehat Q^{\sim}={} \& X_0\arrow[ur]\arrow[dr]\arrow[rr,dotted] \& \& \widetilde{\mathscr X}_{\unsure{\overline Q|Q}}\arrow[r] \& Y\\
\& \& X_2\arrow[ur] \& \&
\end{tikzcd}}

\diagram{\begin{tikzcd}[ampersand replacement=\&,row sep=.3em,column sep=1.4em]
\& \& X_1^\circ\arrow[dr,Rightarrow] \& \& \\
\widehat{Q^\circ}^{\sim}={} \& X_0^\circ\arrow[ur]\arrow[dr]\arrow[rr,dotted] \& \& \widetilde{\mathscr X}_{Q^\circ}\arrow[r] \& Y^\circ\\
\& \& X_2^\circ\arrow[ur,dash] \& \&
\end{tikzcd}}

\deleted{\unsure{déjà|déf.}} sont isomorphes dans $\mathrm{Hot}_W(\Gamma)$. De façon plus précise encore, on a les isomorphismes dans \uline{$\mathrm{Hom}$}$(\Gamma^\circ,\mathrm{Cat})$ :

\[
\begin{aligned}
(\widehat Q^{\sim})^\circ &\simeq \widehat{Q^\circ}\\
(\widehat{Q^\circ}^{\sim})^\circ &\simeq \widehat Q,
\end{aligned}
\]

\margin{d'où en l'appliquant à $Q^\circ$ au lieu de $Q$}

ce qui implique (vu les isom. canoniques valables dans $\mathrm{HOT}_W(\unsure{S|D})$ pour \deleted{toute catégorie} objet $S$ \deleted{dans} $\mathrm{Cat}$ et tout $X_*\in\mathrm{Ob}\underline{\mathrm{Hom}}(S,\mathrm{Cat})$, $X_*\simeq X_*^\circ$${}^{\star}$) qu'on a les isom. canoniques
% slanted
\margin{via le foncteur $\unsure{\Sigma|Z}$ utilisé plus haut\ldots}

\[
\left.\begin{aligned}
\widehat Q&\simeq\widehat{Q^\circ}^{\sim}\\
\widehat{Q^\circ}&\simeq\widehat Q^{\sim}
\end{aligned}\right|\quad\text{dans }\mathrm{HOT}_W(\Gamma)
\]

mais aussi directement

\[
\widehat Q\simeq\widehat{Q^\circ},\quad \widehat Q^{\sim}\simeq\widehat{Q^\circ}^{\sim}\quad\text{dans }\mathrm{HOT}_W(\Gamma)
\]
% archive sheet number: 193
\scanpage{194}\mspage{143}
(et il reste à vérifier qu'on a ainsi un syst. \uline{transitif} d'isom. canoniques dans $\mathrm{HOT}_W(\Gamma)$, entre les quatre objets envisagés, \deleted{\unsure{)}} ce qui doit être essentiellement \unsure{sorital}\ldots).

\uline{Remarques} 1.) Ce qu'on vient de dire dans le cas des diagrammes en $\mathrm{Cat}$ de type $\Phi^\circ=\Psi$, \inserted{\unsure{a|à} l'air de} garder un sens pour \unsure{les|des} diagrammes de type quelconque

\[
X_*:S^\circ\longrightarrow\mathrm{Cat},\qquad s\longmapsto X_s .
\]

($S$ un objet de $\mathrm{Cat}$). On peut l'intégrer de façon lisse, $\int_{\deleted{\unsure{\ell}}}X_*\underset{\unsure{\mathrm{Df}|{D}}}{=}\int_{\mathrm{lis}}X_*$, obtenant une catégorie fibrée (scindée) sur $S$, ou de façon propre, $\int_{\mathrm{prop}}X_*$, obtenant une catégorie cofibrée scindée sur $S^\circ$. On peut également considérer $X_*^\circ$, composé de $X_*$ avec l'autom. $Y\longmapsto Y^\circ$ de $\mathrm{Cat}$, et l'intégrer de façon soit propre, soit lisse. Dans tous ces cas, on trouve un diagramme de type $\widehat S^\circ$, où $\widehat S$ est déduit de $S$ en \unsure{appendant} un élément \deleted{initial} \deleted{\unsure{décrit}} \deleted{\unsure{caractérisé par}} \inserted{qui} \deleted{une flèche \unsure{unique} dans tous les autres (sans} s'envoie par une flèche ``canonique'' \inserted{(\unsure{choisie|donnée|libre?})} dans chacun des autres (\uline{sans} relations de commutation !). Un
% archive sheet number: 194
\scanpage{195}\mspage{144}
\uline{système inductif} de morphismes dans $\mathrm{Cat}$

\[
f\qquad X_s\xrightarrow{f_s}Y
\]

donne donc un morphisme

\[
\int_{\mathrm{lis}}X_*\xrightarrow{\varphi_{\unsure{Q}}}Y,
\]

où $Q$ désigne le diagramme, de type $S^{\wedge\,\circ}$, formé des $X_s$, de $Y$, des $u^*:X_t\to X_s$ (pour $u:s\to t$ dans $S$) et des $f_s$. De $=$ on aura

\[
\int_{\mathrm{prop}}X_*\xrightarrow{\widetilde\varphi_Q}Y,
\]

et si

\[
\Gamma=(S^\wedge+\text{objet initial}),
\]

on trouve des diagrammes $\widehat Q,\widehat Q^{\sim}$ de type $\Gamma^\circ$ dans $\mathrm{Cat}$. Sauf erreur, les relations entre $\widehat Q,\widehat Q^{\sim},\widehat{Q^\circ},\widehat{Q^\circ}^{\sim}$ sont les $\unsure{{=}^{\mathrm s}|{=}\mathrm s}$ que celles explicitées dans la proposition 14 ci-dessus. Je ne vais pas m'arrêter là-dessus avec \unsure{le soin requis|la sur-rigueur}, sauf en cas de besoin. On dira qu'un diagramme $Q$ définit $Y$ comme $\underset{\unsure{I|1}}{\varinjlim}$ $W$-homotopique, si $\varphi_Q\in W$\deleted{\unsure{,} \illegible{2} \unsure{de} \unsure{$f$}} Il revient au $=$ de dire que $\varphi_{Q^\circ}\in W$, ou que $\widetilde\varphi_Q\in W$, ou que $\widetilde\varphi_{Q^\circ}\in W$\ldots.
% archive sheet number: 195
\scanpage{196}\mspage{145}
2) Le cas ``universel'', pour

\diagram{\begin{tikzcd}[ampersand replacement=\&,row sep=0pt,column sep=1.2em,cramped] \& X_1 \\ X_0\arrow[ur]\arrow[dr] \& \\ \& X_2 \end{tikzcd}}

donné, d'un carré co\deleted{\unsure{$\varphi_Q$}}mmutatif

\diagram{\(Q={}\begin{tikzcd}[ampersand replacement=\&,row sep=0pt,column sep=1.4em,cramped] \& X_1\arrow[dr] \& \\ X_0\arrow[ur]\arrow[dr] \& \& Y \\ \& X_2\arrow[ur] \& \end{tikzcd}\)}

dans $\mathrm{Cat}$, donnant naissance à

\diagram{\(\varphi_Q:\left(\int \begin{tikzcd}[ampersand replacement=\&,row sep=0pt,column sep=1.2em,cramped] \& X_1 \\ X_0\arrow[ur,dash]\arrow[dr] \& \\ \& X_2 \end{tikzcd}\right)\longrightarrow Y,\)}

est celui de $Y=X_{\unsure{1|0}}\amalg_{X_0}X_2$, donnant naissance à

\diagram{\(Q_0=Q(X_*):\quad \begin{tikzcd}[ampersand replacement=\&,row sep=0pt,column sep=1.4em,cramped] \& X_1\arrow[dr] \& \\ X_0\arrow[ur]\arrow[dr] \& \& X_1\amalg_{X_0}X_2 \\ \& X_2\arrow[ur] \& \end{tikzcd}\,.\)}

Le cas général s'y ramène car on aura un diagramme commutatif

\diagram{\begin{tikzcd}[ampersand replacement=\&,column sep=4em,row sep=2.5em]
{\left(\int \begin{tikzcd}[ampersand replacement=\&,row sep=0pt,column sep=1.2em,cramped] \& X_1 \\ X_0\arrow[ur]\arrow[dr] \& \\ \& X_2 \end{tikzcd}\right)}\arrow[r,"{\varphi_{Q_0}=\varphi_{Q(X_*)}}"]\arrow[dr,"\varphi_Q"'] \& X_{\unsure{4|1|t}}\amalg_{X_0}X_2\arrow[d,"\alpha_Q"]\\
\& Y
\end{tikzcd}}

Mais il doit y avoir des cas intéressants où $\varphi_{Q_0}$ n'est \uline{pas} dans $W$, i.e. le carré \inserted{$\mathrm{Cat}$-}cocartésien

\diagram{\begin{tikzcd}[ampersand replacement=\&,row sep=0pt,column sep=1.4em,cramped] \& X_1\arrow[dr] \& \\ X_0\arrow[ur]\arrow[dr] \& \& X_1\amalg_{X_0}X_2 \\ \& X_2\arrow[ur] \& \end{tikzcd}}

n'est pas $W$-homotopiquement cocartésien, et où pourtant on peut trouver un carré commutatif $Q$
% archive sheet number: 196
\scanpage{197}\mspage{146}
complétant $X_*$, de telle façon que ce carré soit $W$-homotopiquement commutatif, i.e que $\varphi_Q=\alpha_Q\varphi_{Q_0}$ soit dans $W$. C'est d'autant plus plausible que dans $\mathrm{Cat}$, les sommes amalgamées ont \unsure{tendance} à ne pas avoir des bonnes propriétés homotopiques, même si l'une des flèches $\alpha_1,\alpha_2$ est mono, (voire une immersion ouverte), contrairement à ce qui se passe dans $\Delta^\wedge$, et plus généralement dans les catégories de la forme $Z^\wedge$ ($Z$ dans $\mathrm{Cat}$).

3.4. \uline{Critères pour les carrés homotopiquement cocartésiens}.

Les critères que j'ai dégagés sont tous contenus dans l'énoncé suivant :

\uline{\deleted{Proposition \unsure{11}}}\inserted{Théorème 4} Considérons un \inserted{diagramme} carré commutatif \unsure{dans|de} $\mathrm{Cat}$

% slanted
\margin{\uline{NB} Dans ce th., et les cor. 1, 2, 3, on n'a besoin que de la \uline{\unsure{forme} faible} de $W(7\text{bis})$, savoir que $\unsure{ir|i\circ r|i_0}\in W\Longrightarrow i\in W$ (cf p. \unsure{74|76}), (Loc($6'\mathrm a$) avec les notations de XVI)}

% slanted
\margin{\uline{NB} \unsure{Au total|En fait}, ce théorème \inserted{(\unsure{partie directe|partie droite})} est valable sans hypothèse $W(7\text{bis})$ ou $W(7\text{bis }\unsure{\text{faible}|\text{fort}})$ sur $W$, et il en est de $=$ de ses corollaires 1 à 3 (à l'exclusion des corollaires 4, 5), qui eux utilisent $W(7\text{bis})$ dans \unsure{tte} sa force).}

\diagram{\(Q={}\begin{tikzcd}[ampersand replacement=\&,row sep=1.6em,column sep=3em]
X_0\arrow[r,"\alpha_1"]\arrow[d,"\alpha_2"']\arrow[dr,dashed,"\unsure{f_0}"] \& X_1\arrow[d,"f_1"]\\
X_2\arrow[r,"f_2"'] \& Y
\end{tikzcd}\)}

et considérons le morphisme canonique

\diagram{\(\varphi_Q:\unsure{\mathscr X|\mathcal X}=\left(\int \begin{tikzcd}[ampersand replacement=\&,row sep=0pt,column sep=1.2em,cramped] \& X_1 \\ X_0\arrow[ur,"\alpha_1"]\arrow[dr,"\alpha_2"'] \& \\ \& X_2 \end{tikzcd}\right)\longrightarrow Y,\)}
% archive sheet number: 197
\scanpage{198}\mspage{147}
l'intégrale étant l'intégrale \uline{lisse}. (Il y a un énoncé dual concernant l'intégrale propre de

\diagram{\begin{tikzcd}[ampersand replacement=\&,row sep=0pt,column sep=1.2em,cramped] \& X_1 \\ X_0\arrow[ur]\arrow[dr] \& \\ \& X_2 \end{tikzcd}\(\,.\))}

\deleted{Pour ce} Pour que $\varphi_Q$ soit \inserted{$W$-}lisse, il faut et s. que $f_1,f_2,\unsure{f_0}=f_1\alpha_1=f_2\alpha_2$ le soient. (\uline{Moyennant $W(7\text{bis})$})

% slanted
\margin{Pour la partie directe ($f_0,f_1,f_2$ lisses $\Longrightarrow\varphi_Q$ lisse), on n'a besoin que de $W(7\text{bis})$ forme faible. Ce n'est que pour la réciproque (qu'on n'aura jamais à utiliser) qu'on a besoin de $W(7\text{bis})$ (et dans tte sa force).}

% slanted
\margin{On n'a pas $=$ besoin de $W(7\text{bis})$ \inserted{faible}, cf annotation page 156. On utilise une ``\unsure{hypothèse}'' \unsure{de lissité|de rétraction} sur \unsure{$W$|les} \unsure{pour} la \unsure{$\Longrightarrow$} \unsure{réciproque} seulement ($\varphi$ lisse $\Longrightarrow f_0,f_1,f_2$ lisses).}

La fibre de $\varphi_Q$ en un \deleted{\unsure{pt}} élément \unsure{$y$} est isomorphe à

\diagram{\(\text{(**)}\qquad\deleted{\unsure{y}}\qquad \int \begin{tikzcd}[ampersand replacement=\&,row sep=0pt,column sep=1.3em,cramped] \& X_{1y}\\ X_{0y}\arrow[ur]\arrow[dr] \& \\ \& X_{2y}\end{tikzcd}\quad (\simeq\unsure{\mathscr X|\mathcal X}_y).\)}

\uline{Corollaire 1} Supposons $f_1,f_2$ et $f_1\alpha_1=f_2\alpha_2$ $W$-lisses, et les catégories intégrales $(**)$ \deleted{\unsure{de}} \inserted{$W$-}asphériques pour \deleted{les} $y\in\mathrm{Ob}\,Y$. [NB Il s'agit ici de l'intégrale lisse, mais peu importe en fait, car on a vu \deleted{que} (cf \deleted{volet} \deleted{\unsure{gauche}} partie supérieure des diagrammes p. 140, qui impliquent que $\mathscr X$ $W$-asph. \unsure{sss|ssi} $\widetilde X$ l'est) que l'une est $W$-asph. sss l'autre l'est.]

% slanted
\margin{Valable sans hypothèse sur $W$ (autre que Loc(1,2,\unsure{3bis})).}

Alors $\varphi_Q\in W$, i.e. le carré $\unsure{(\alpha)|(*)}$ est $W$-homotopiquement \deleted{trivial}\inserted{cartésien}. Par dualité, $=$ conclusion si on suppose $f_1,f_2,f_1\alpha_1=f_2\alpha_2$ \uline{propres} au lieu de lisses. (car $Q$ $W$-homotopiquement trivial, sss $Q^\circ$ l'est, par prop. 14.)
% archive sheet number: 198
\scanpage{199}\mspage{148}
En effet, un morphisme dans $\mathrm{Cat}$ qui est $W$-lisse et à fibres $W$-asphériques est dans $W$ (et $=$ \inserted{c'est une $W$-équivalence} universelle \deleted{\unsure{ment}}\ldots)

\uline{Corollaire 2} \inserted{On suppose que $W$ satisfait $W(7\text{bis})$ (axiome du carré cocartésien).}\margin{\unsure{inutile}} \deleted{Considérons un}\inserted{Supposons le} carré $Q$ cocartésien dans $\mathrm{Cat}$${}^{*}$, et que $\alpha_1$ soit une immersion ouverte (resp. fermée) et $f_1$ lisse (resp. propre). Alors $Q$ est \inserted{$W$-}homotopiquement cocartésien.

% slanted
\margin{NB On n'utilise \unsure{pas|\deleted{pas}} vraiment que $Q$ soit cocartésien, plutôt que $Q$ soit \uline{cartésien}, et $f_2$ (donc aussi $\alpha_1$) une imm. ouverte (resp. fermée), et que pour $y\in Y\setminus X_2$, $X_{1y}$ soit $W$-asphérique. (Dans le cas où $Q$ est cocartésien $X_{1y}$ est $\simeq e$.)}

\margin{On a besoin ici de $W(7\text{bis})$ \uline{forme directe}, i.e. de $\mathrm{Loc}(6'\mathrm a)$ (notation de \uline{XVI})}

Il suffit d'établir le résultat non respé, l'autre en résulte par dualité (puisque $Q$ est $W$-\deleted{homotopiquement}\inserted{cocartésien} sss $Q^\circ$ l'est).

Par hypothèse $f_1$ est \inserted{$W$-}lisse, d'autre part $\alpha_1$ et $f_2$ sont des imm. ouvertes donc $W$-lisses, donc $f_1\alpha_1$ est \inserted{$W$-}lisse aussi. (D'ailleurs $\alpha_2$ aussi, car le carré est aussi cartésien donc $\alpha_2$ déduit de $f_1$ \inserted{(qui est $W$-lisse par hypothèse)} par chgt. de base $X_2\to Y$.) Donc $\unsure{\varphi|f}_Q$ est lisse, et il suffit de voir que ses fibres sont $W$-asphériques. Le calcul \unsure{les} donne, pour $y$ dans $Y\setminus X_2$

\diagram{\(\unsure{\mathscr X|\mathcal X}_y\simeq\int\deleted{\unsure{\varnothing}} \begin{tikzcd}[ampersand replacement=\&,row sep=0pt,column sep=1.2em,cramped] \& e \\ \varnothing\arrow[ur]\arrow[dr] \& \\ \& \varnothing \end{tikzcd}\simeq e,\quad\text{OK.}\)}

Et pour $y$ dans $X_2$

\diagram{\(\mathscr X_{\unsure{y|\smash{\overset{\circ}{y}}}}\simeq\int \begin{tikzcd}[ampersand replacement=\&,row sep=0pt,column sep=1.2em,cramped] \& X_{1y} \\ X_{0y}\arrow[ur,"\sim"]\arrow[dr] \& \\ \& e \end{tikzcd},\)}
% archive sheet number: 199
\scanpage{200}\mspage{149}
Comme $X_{0y}\longrightarrow X_{1y}$ est un iso donc est dans $W$, s'ensuit que $e\to\unsure{\mathscr X|\mathcal X}_y$ l'est aussi (prop. 13 p. 125), moyennant $W(7\mathrm{bis})$, donc $\mathscr X_y$ est $W$-asphérique, OK.

% slanted
\margin{On n'a besoin ici que de la forme directe de $W(7\mathrm{bis})$, i.e (notations de \uline{XVI}) de $\mathrm{Loc}(6'\mathrm a)$.}

\uline{Corollaire 3} \inserted{(Moyennant $W(7\mathrm{bis})$)}\margin{inutile} Soit $X$ dans $\mathrm{Cat}$, $U_1,U_2$ deux ouverts de $X$ tels que $X=U_1\cup U_2$, $U_0=U_1\cap U_2$ l'ouvert intersection. Alors le carré d'imm. ouvertes

\diagram{\begin{tikzcd}[ampersand replacement=\&,row sep=1.5em,column sep=2em]
U_0\arrow[r,hook]\arrow[d,hook] \& U_1\arrow[d,hook]\\
U_2\arrow[r,hook] \& X
\end{tikzcd}\(\,.\)}

est $W$-cocartésien. (\unsure{Énoncé} dual pour deux fermés $T_1,T_2$ tels que $T_1\cup T_2=\unsure{x|X}$, avec $T_0=T_1\cap T_2$\ldots)

On applique le cor. 2. On sait que le carré est cocartésien dans $\mathrm{Cat}$${}^{*}$, de plus $U_1\hookrightarrow X$ est bien $W$-lisse, \deleted{\unsure{étant}} une imm. ouverte.

% slanted
\margin{On n'a pas même besoin de le savoir, cf annotation du cor. 2}

\uline{Corollaire 4} \inserted{(Moyennant $W(7\mathrm{bis})^{*}$)} Soit $U\overset{i}{\hookrightarrow}X$ une imm. ouverte telle que $T=X-U$ soit isom. à la catégorie ponctuelle $e$. Pour que $i\in W$, il f. et s. que $i$ soit $W$-asphérique, i.e. que $U_{/t}$ \deleted{soit} (où $t$ est l'unique objet de $T$) soit $W$-asphérique.

% slanted
\margin{Ici on \unsure{a} besoin de $W(7\mathrm{bis})$ forme forte i.e. $i\circ r\in W\unsure{\Longleftarrow|\Longleftrightarrow}i\in W$, cf annotation Th. 4.}

Il suffit bien sûr de prouver que \deleted{\unsure{$i$}}

\[
i\in W\Longrightarrow U_{/t}\ \text{$W$-asphérique}.
\]
% archive sheet number: 200
\scanpage{201}\mspage{150}
Considérons le carré cartésien

\diagram{\begin{tikzcd}[ampersand replacement=\&,row sep=1.2em,column sep=2em]
U_{/t}\arrow[r,hook,"i_{/t}"]\arrow[d,"p"'] \& X_{/t}\arrow[d,"q"]\\
U\arrow[r,hook,"i"'] \& X
\end{tikzcd}\(\,,\)}

nous allons prouver qu'il est $W$-cocartésien. Moyennant $W(7\mathrm{bis})$, il s'ensuivra que $i\in W$ implique $i_{/t}\in W$, et comme $X_{/t}$ est \inserted{$W$-}asphérique (ayant un objet final), cela signifie que $U_{/t}$ est $W$-asphérique, ce qu'on veut prouver.

Il suffit, en vertu du cor. 2, de prouver que le carré est cocartésien dans $\mathrm{Cat}$, puisque $i_{/t}$ est une imm. ouverte ($i$ l'étant), et $X_{/t}\to X$ est un iso. local donc \inserted{$W$-}lisse. On voit que $X_{/t}\setminus U_{/t}\simeq e$ (catégorie \deleted{\unsure{finie|finale}} \inserted{ponctuelle}), \deleted{\unsure{donc $X_{/t}$ déduit de $U_{/t}$ en}} \inserted{d'ailleurs l'unique \unsure{objet|élément} de $X_{/t}\setminus U_{/t}$} est l'élém. final $t_0:t\xrightarrow{\mathrm{id}_t}t$ de $X_{/t}$. Ainsi $H_{t_0}\simeq e_{U_{/t}}$ (préfaisceau final), et $p_!(H_{t_0})\simeq \unsure{\mathbb H|H}'_t\overset{\mathrm{déf}}{=}(u\longmapsto\mathrm{Hom}_X(u,t))$, \deleted{\unsure{donc}} ce qui caractérise la somme amalgamée dans $\mathrm{Cat}$\ldots.

Le cor. 4 n'est autre que le corollaire non démontré de la prop \unsure{7} (p. \unsure{85}).

Voici une variante du cor. 4

\uline{Coro. 5} \inserted{(Moyennant $W(7\mathrm{bis})^{*}$)}
% slanted
\margin{forme forte.} Soit $i:U\hookrightarrow X$ une imm. ouverte, soit $t_0\in\mathrm{Ob}\,T$ un point \uline{fermé} de $T=X\setminus U$, tel que $\{t_0\}$
% archive sheet number: 201
\scanpage{202}\mspage{151}
soit iso. à la catégorie ponctuelle, i.e $\mathrm{Hom}_X(t_0,t_0)=\{\mathrm{id}_{\unsure{t|t_0}}\}$. Supposons que pour tout $t\in T\setminus\{t_0\}$, $U_{/t}$ soit $W$-asphérique. Pour que l'on ait $i\in W$, il f et s. que $i$ soit asphérique, i.e que $U_{/t_0}$ soit également $W$-asphérique.

Il faut prouver que $i\in W$ implique que $U_{/t_0}$ est $W$-asphérique. Considérons l'inclusion

\[
i^{\prime}_0:U\longrightarrow \deleted{\unsure{X}}U'\overset{\mathrm{déf}}{=}X\setminus\{t_0\},
\]

l'hypothèse sur les $U_{/t}$ \inserted{pour $t\in T-\{t_0\}$} signifie que $i'$ est $W$-asphérique, donc $i'$ est dans $W$, et comme \deleted{\unsure{le}} le composé $i:U\xrightarrow{\unsure{i'_0|i_0}}U'\xrightarrow{i'}X$ est dans $W$, cela implique que $i':U'\mathbin{\unsure{-|\to}}X$ est dans $W$. Nous sommes maintenant dans la situation du cor. 4, qui implique que $i'$ est $W$-asphérique. Donc $i=i'i'_0$ est $W$-asphérique, comme composé de deux morphismes $W$-asphériques, qed.

\uline{Démonstration du \deleted{prop. \unsure{15|11}} \inserted{th. 4}} \textcircled{a} Critère de lissité de

\diagram{\(\varphi_Q:\unsure{\mathscr X|\mathcal X}=\int \begin{tikzcd}[ampersand replacement=\&,row sep=0pt,column sep=1.2em,cramped] \& X_1 \\ X_0\arrow[ur,dash]\arrow[dr] \& \\ \& X_2 \end{tikzcd}\longrightarrow Y.\)}

Comme $\{1\},\{\unsure{1|2}\}$ sont des \deleted{\unsure{sous-cat.}} ouverts de

\diagram{\(\Phi={}\begin{tikzcd}[ampersand replacement=\&,row sep=0pt,column sep=1.4em]
\& 1\arrow[dl]\\0\&\\\&2\arrow[ul]\end{tikzcd},\)}

\deleted{\unsure{$X_1$}}, $i_1:X_1\to\mathscr X$ et $i_2:X_2\to\mathscr X$ sont des imm. ouvertes. Comme $\varphi_Q\circ i_{\unsure{\nu|v|j}}=f_\nu$,
% archive sheet number: 202
\scanpage{203}\mspage{152}
$\varphi_Q\circ i_2=f_{\unsure{2|\nu}}$, on voit déjà que

\[
\begin{aligned}
f_1\text{ et }f_2\text{ lisses}\Longleftrightarrow {}&\varphi_Q\text{ lisse sur l'ouvert réunion}\\
&\text{de }i_1(X_1)\text{ et }i_2(X_2),\ \text{i.e.}\\
&\text{sur }\unsure{\mathscr X|\mathcal X}\setminus i_0(X_0).
\end{aligned}
\]

Il reste à exprimer la lissité de $\varphi_Q$ en les points $x_0$ de $i_0(X_0)$. Soient $x_1,x_2,y$ les images de \deleted{\unsure{$x_0$ point}} dans $X_1,X_2,Y$ respectivement. Il faut exprimer que les fibres de \unsure{les} $(\varphi_Q)_{/x_0}\quad\mathscr X_{/x_0}\longrightarrow Y_{/y}$ sont asphériques. \deleted{Mais} Nous allons pour cela expliciter $\mathscr X_{/x_0}$ \inserted{et la flèche \unsure{can|sus}}, en termes du carré commutatif

\diagram{\((**)\qquad\begin{tikzcd}[ampersand replacement=\&,row sep=1.7em,column sep=3em]
X_{0/x_0}\arrow[r,"{\alpha_{1/x_0}}"]\arrow[d,"{\alpha_{2/x_0}}"'] \& X_{1/x_1}\arrow[d,"{f_{1/x_1}}"]\\
X_{2/x_2}\arrow[r,"{f_{2/x_2}}"] \& Y_{/y}
\end{tikzcd}\)}

induit par $Q$. Je dis

\uuline{Lemme 1} \unsure{Soit|Soient} \deleted{\unsure{$x_0$}}un carré commutatif \deleted{$Q$}

\diagram{\(\unsure{\in|\epsilon}\qquad Q=\deleted{\unsure{Q\to}}\begin{tikzcd}[ampersand replacement=\&,row sep=1em,column sep=2em]
X_0\arrow[r,"\alpha_1"]\arrow[d,"\alpha_2"'] \& X_1\arrow[d,dash,"f_1"]\\
X_2\arrow[r,"f_2"] \& Y
\end{tikzcd}\)}

dans $\mathrm{Cat}$, $x_0\in\mathrm{Ob}\,X_0$, $x_1,x_2,y$ ses images dans $X_1,X_2,Y$ respectivement, $Q_{/x_0}$ le carré \unsure{$(**)$|corresp.} ci-dessus, induit par $Q$ sur les localisés $X_{0/x_0}$ etc.,

\diagram{\(\mathscr X=\int \begin{tikzcd}[ampersand replacement=\&,row sep=0pt,column sep=1.2em,cramped] \& X_1 \\ X_0\arrow[ur]\arrow[dr,dash] \& \\ \& X_2 \end{tikzcd}.\)}

On a alors un
% archive sheet number: 203
\scanpage{204}\mspage{153}
isomorphisme canonique

\diagram{\(\unsure{\mathscr X|\mathcal X}_{/x_0}\simeq\int \begin{tikzcd}[ampersand replacement=\&,row sep=0pt,column sep=1.2em,cramped] \& X_{1/x_1} \\ X_{0/x_0}\arrow[ur]\arrow[dr] \& \\ \& X_{2/x_2} \end{tikzcd},\)}

et la flèche $(\varphi_Q)_{/x_0}$ localisée de $\varphi_Q:\mathscr X\mathbin{\unsure{-|\to}}Y$ en \deleted{\unsure{$y$}} $x_0$ s'identifie à la flèche $\varphi(Q_{/x_0})$ :

\diagram{\(\varphi(Q_{/x_0}):\ \int \begin{tikzcd}[ampersand replacement=\&,row sep=0pt,column sep=1.2em,cramped] \& X_{1/x_1} \\ X_{0/x_0}\arrow[ur]\arrow[dr] \& \\ \& X_{2/x_2} \end{tikzcd}\longrightarrow Y_{/y},\)}

\uline{Démonstration} \inserted{du lemme.} On sait que

\[
\mathscr X_{/x_0}\longrightarrow\mathscr X
\]

est \inserted{$\mathrm{Cat}$-}fibrée \inserted{(scindée)} (à fibres discrètes), et comme

\[
\mathscr X\longrightarrow\Phi
\]

est aussi $\mathrm{Cat}$-fibrée (scindée), il s'ensuit qu'il en est de $=$ de la flèche composée

\[
\mathscr X_{/x_0}\longrightarrow\Phi .
\]

Donc \deleted{il faut} \unsure{,} pour ``calculer'' $\mathscr X_{/x_0}$, en \deleted{\unsure{des}} que catégorie sur $\Phi$, il faut expliciter ses \inserted{trois} fibres, et les deux foncteurs de chgt. de base qui les relient. On trouve aussitôt que ces fibres sont can. isomorphes à

\[
X_{0/x_0},\quad X_{1/x_1},\quad X_{2/x_2},
\]

et on admettra (!) que les foncteurs chgt de
% archive sheet number: 204
\scanpage{205}\mspage{154}
\[
X_{0/x_0}\longrightarrow X_{1/x_1},\qquad X_{0/x_0}\mathbin{\unsure{-|\to}} X_{2/x_2}
\]

\inserted{(définis à iso. uniques près)} ne sont autres que \unsure{les|\deleted{les}} localisées $\alpha_{1/x_0}$ et $\alpha_{2/x_0}$ de $\alpha_1:X_0\mathbin{-}X_1$ et $\alpha_2:X_0\mathbin{-}X_2$. Cela montre donc qu'on a un iso.

\diagram{\(\unsure{\mathscr X|\mathcal X}_{/x_0}\simeq\int \begin{tikzcd}[ampersand replacement=\&,row sep=0pt,column sep=1.2em,cramped] \& X_{1/x_1} \\ X_{0/x_0}\arrow[ur]\arrow[dr] \& \\ \& X_{2/x_2} \end{tikzcd},\)}

et l'on admettra aussi que la flèche $(\varphi_Q)_{/x_0}:\mathscr X_{/x_0}\mathbin{-}Y_{/y}$ s'identifie à

\[
\varphi(Q_{/x_0}).
\]

\unsure{Pour} exprimer la \inserted{$W$-}lissité de $\varphi_Q$ en $x_0$\inserted{$\in\mathrm{Ob}\,X_0\subset\mathrm{Ob}\,\mathscr X$}, on est donc ramené à exprimer la $W$-asphéricité des fibres de

\diagram{\(\varphi(Q_{/x_0}):\ \int \begin{tikzcd}[ampersand replacement=\&,row sep=0pt,column sep=1.2em,cramped] \& X_{1/x_1} \\ X_{0/x_0}\arrow[ur]\arrow[dr] \& \\ \& X_{2/x_2} \end{tikzcd}\longrightarrow Y_{/y}.\)}

Pour expliciter les fibres, on a besoin d'un autre lemme auxiliaire

\uline{Lemme 2} $^{\star}$ Soit un carré commutatif dans $\mathrm{Cat}$

\margin{NB \unsure{C'est le} b) \unsure{du Th} 4 \deleted{\unsure{page 148}}} % slanted

\diagram{\(Q=\biggl(\begin{tikzcd}[ampersand replacement=\&,row sep=1.0em,column sep=1.8em] X_0\arrow[r,"\alpha_1"]\arrow[d,"\alpha_2"'] \& X_1\arrow[d,"f_1"] \\ X_2\arrow[r,"\deleted{\unsure{f_0}}" above,"f_2" below] \& Y \end{tikzcd}\biggr),\)}

et \deleted{\unsure{soit}}

\diagram{\(\varphi_Q:\bigl(\int \begin{tikzcd}[ampersand replacement=\&,row sep=0pt,column sep=1.2em,cramped] \& X_1 \\ X_0\arrow[ur]\arrow[dr,dash] \& \\ \& X_2 \end{tikzcd}\bigr)\longrightarrow Y,\)}
% archive sheet number: 205
\scanpage{206}\mspage{155}
et soit $y\in\mathrm{Ob}\,Y$. Alors la fibre de $\varphi_Q$ en $y$

\[
\unsure{\mathscr X|\mathcal X}_y=\varphi_Q^{-1}(\underset{\substack{\rule{.35pt}{.4ex}\rule{1.6em}{.35pt}\rule{.35pt}{.4ex}\\[-.2ex]\text{cat. ponctuelle}}}{\{y\}})
\]

est donnée \deleted{\unsure{par}} (comme cat. au dessus de $\Phi$)

\diagram{\(\mathscr X_y\simeq\int\begin{tikzcd}[ampersand replacement=\&,row sep=0pt,column sep=1.2em,cramped] \deleted{\unsure{X_{0y}}}\arrow[r] \& X_{1y}\\ X_{0y}\arrow[ur]\arrow[dr] \& \\ \& X_{2y} \end{tikzcd}\)}

\uline{Démonstration} : Il est immédiat qu'on a

\[
\mathscr X_y\cap\underset{\substack{\rule{.35pt}{.4ex}\rule{1.4em}{.35pt}\rule{.35pt}{.4ex}\\[-.2ex]X_s}}{\mathscr X_s}\simeq(X_s)_y\qquad\text{pour }\ s\in\mathrm{Ob}\,\Phi=\{0,1,2\},
\]

\margin{\diagram{\begin{tikzcd}[ampersand replacement=\&,row sep=1.1em,column sep=1em,cramped] X_{sy}\arrow[r]\arrow[d] \& \mathscr X_y\arrow[r]\arrow[d]\arrow[dl,phantom,"\mathrm{cart}"] \& \{y\}\arrow[d]\arrow[dl,phantom,"\mathrm{cart}"] \\ X_s\arrow[r]\arrow[d] \& \mathscr X\arrow[r]\arrow[d]\arrow[dl,phantom,"\mathrm{cart}"] \& Y \\ \{s\}\arrow[r] \& \Phi \& \end{tikzcd}}} % left margin beside the intersection

et (le hasard fait bien les choses) les foncteurs chgt. de base

\[
\begin{aligned}
&\deleted{X_0\longrightarrow X_1}\\
&X_0\longrightarrow X_1,\qquad X_0\longrightarrow X_2
\end{aligned}
\]

\deleted{\unsure{transforment}} \inserted{envoient} la sous-catégorie (non pleine) $X_{0y}$ de $X_0$ dans \inserted{les sous-catégories (non pleines)} $X_{1y}$ et \deleted{\unsure{dans}} $X_{2y}$ de $X_1$ et $X_2$ respectivement. On a donc une sous-catégorie fibrée de $\mathscr X$, sous réserve de vérifier que les flèches \unsure{strictmt|strictt|strictement} cartésiennes $x_0\longrightarrow\alpha_1(x_0)$, $x_0\mathbin{\unsure{-|\to}}\alpha_2(x_0)$ dans $\mathscr X$, \deleted{\unsure{sont}} pour $x_0$ dans $X_{0y}$, sont dans $\mathscr X_y$, i.e sont au dessus de la flèche identique de $y$. Or c'est là la propriété essentielle
% archive sheet number: 206
